\begin{abstract}
The intricate and highly structured characteristics of mathematical reasoning create substantial difficulties for language models. This work presents DeepSeekMath~7B, which extends the pre-training of DeepSeek-Coder-Base-v1.5 7B using 120B tokens related to mathematics, collected from Common Crawl alongside additional natural language and programming data. Notably, DeepSeekMath~7B attains a notable 51.7\% on the MATH benchmark, a competition-level evaluation, all without the use of external support tools or ensembling methods. This performance narrows the gap with leading proprietary models like Gemini-Ultra and GPT-4. Furthermore, sampling 64 responses from DeepSeekMath~7B and aggregating results through self-consistency yields a 60.9\% score on MATH. The advancements in mathematical reasoning observed in DeepSeekMath~stem from two principal innovations: firstly, an extensive and carefully constructed data curation process that maximally leverages freely accessible web resources; secondly, the development of Group Relative Policy Optimization (GRPO)—an adaptation of Proximal Policy Optimization (PPO)—which both strengthens mathematical reasoning capabilities and improves PPO's memory efficiency.
\end{abstract}

\section{Introduction}

The advent of large language models (LLMs) has transformed artificial intelligence methods for mathematical reasoning, catalyzing substantial progress on tasks such as quantitative reasoning benchmarks~\citep{MATH} and geometric problem solving~\citep{trinh2024solving}. Beyond benchmark results, these systems are increasingly valuable in supporting human users as they tackle advanced mathematical questions~\citep{tao}. Despite these advancements, top-performing models like GPT-4~\citep{gpt4} and Gemini-Ultra~\citep{gemini} remain closed-source, leaving available open-source alternatives far behind in terms of capability.

This work presents DeepSeekMath, a specialized language model tailored for mathematics that surpasses other open-source alternatives and attains performance nearing that of GPT-4 on established academic evaluations. To facilitate this advancement, we constructed the DeepSeekMath Corpus—an extensive, high-quality pre-training dataset comprising 120B mathematical tokens. This corpus was assembled from Common Crawl (CC) resources using a classifier built with fastText \citep{joulin2016fasttext}. For its initial training phase, the classifier utilized positive examples from OpenWebMath \citep{openwebmath}, complemented by a diverse array of unrelated web content serving as negative samples. Afterward, the classifier was deployed to retrieve more positive samples from the CC, and these selections were further validated and improved by human annotators. The enhanced labeled dataset was then used to retrain the classifier, yielding improved extraction capabilities. Empirical assessment demonstrates the strength of our corpus: DeepSeekMath-Base 7B achieves 64.2\% on GSM8K \citep{gsm8k} and 36.2\% on the MATH competition dataset \citep{MATH}, exceeding the results of Minerva 540B \citep{minerva}. The DeepSeekMath Corpus is also multilingual, leading to marked gains on Chinese mathematical benchmarks \citep{wei2023cmath,agieval}. We propose that our approach to processing mathematical data provides a valuable foundation for ongoing research and anticipate substantial future advancements in this area.

DeepSeekMath-Base is built upon DeepSeek-Coder-Base-v1.5 7B \citep{deepseek-coder}, motivated by findings that initializing with a code-oriented model outperforms starting from a general-purpose large language model. Notably, incorporating mathematics-centric training not only augments the model’s proficiency in math, but also yields gains on benchmarks such as MMLU \citep{mmlu} and BBH \citep{bbh}, revealing improvements in general reasoning abilities as well.

Subsequent to its pre-training stage, DeepSeekMath-Base undergoes mathematical instruction tuning, which leverages data from chain-of-thought \citep{cot}, program-of-thought \citep{pot,pal}, and tool-integrated reasoning paradigms \citep{tora}. The fine-tuned DeepSeekMath-Instruct 7B consequently surpasses other 7B models and demonstrates performance on par with the leading 70B open-source models tuned for instructions.

Additionally, we propose Group Relative Policy Optimization (GRPO), a new reinforcement learning algorithm adapted from Proximal Policy Optimization (PPO) \citep{schulman2017proximal}. In contrast to PPO, GRPO dispenses with the critic network, estimating baselines from aggregated group scores and thereby reducing the computational demands of training. Utilizing only a portion of English-language instruction tuning data, GRPO delivers marked performance improvements over the strong DeepSeekMath-Instruct baseline on both in-domain datasets (GSM8K: 82.9\% $\rightarrow$ 88.2\%, MATH: 46.8\% $\rightarrow$ 51.7\%) and out-of-domain math challenges (e.g., CMATH: 84.6\% $\rightarrow$ 88.8\%) during reinforcement learning.

Furthermore, we present a unified conceptual framework for interpreting various approaches such as Rejection Sampling Fine-Tuning (RFT) \citep{yuan2023scaling}, Direct Preference Optimization (DPO) \citep{dpo}, PPO, and GRPO, showing that all can be understood as variations of direct or simplified reinforcement learning methodologies. Through comprehensive experimentation—contrasting online versus offline learning, outcome versus process supervision, and single-step versus iterative RL—we dissect the pivotal aspects of this paradigm. Finally, we elucidate the mechanisms by which reinforcement learning enhances instruction-tuned models, and propose future avenues for advancing RL effectiveness within this unified perspective.

\subsection{Contributions}
This work presents advancements in large-scale mathematical pre-training and delivers a comprehensive examination of reinforcement learning strategies.

\textbf{Math Pre-Training at Scale}

\begin{itemize}[topsep=0pt]
    \item We demonstrate that the widely available Common Crawl corpus is a rich resource for mathematics-related material. Leveraging a rigorously engineered selection process, we assemble the DeepSeekMath Corpus, consisting of 120B tokens drawn from mathematically relevant web pages. This corpus notably surpasses the size of the math web pages in Minerva \citep{minerva} by nearly sevenfold and exceeds the volume of OpenWebMath \citep{openwebmath} by a factor of nine.

    \item Our DeepSeekMath-Base 7B base model, after pre-training, attains results on par with the substantially larger Minerva 540B \citep{minerva}. This observation suggests that, for mathematical reasoning, model size is not the sole determinant of effectiveness; a more compact model pre-trained on curated, high-quality data can perform exceptionally well.

    \item We report insights gleaned from our experiments involving mathematics training. Notably, pre-training with code prior to mathematical datasets enhances the model’s problem-solving abilities in both tool-assisted and tool-free scenarios. This provides empirical support to the question: \textit{does code pre-training strengthen reasoning?} In the context of mathematics, our findings indicate that it does.

    \item Contrary to common practice, our experiments show that pre-training on arXiv papers, although frequently adopted in mathematics-related studies, yields negligible benefits across all the mathematical benchmarks considered herein.
\end{itemize}

\textbf{Exploration and Analysis of Reinforcement Learning}

\begin{itemize}[topsep=0pt]
    \item We propose Group Relative Policy Optimization (GRPO), a novel reinforcement learning algorithm characterized by both efficiency and effectiveness. Unlike traditional methods that rely on a critic network, GRPO derives its baseline directly from group-level scoring, thereby substantially decreasing computational and training costs in comparison to Proximal Policy Optimization (PPO).
    \item Our results show that applying GRPO to instruction-tuned models, specifically DeepSeekMath-Instruct, results in significant improvements by leveraging only instruction-tuning datasets. Additionally, reinforcement learning with GRPO enhances generalization to out-of-domain tasks.
    \item We establish a unified conceptual framework encompassing various techniques such as RFT, DPO, PPO, and GRPO. Through extensive experimentation—including comparisons between online and offline learning, outcome versus process-based supervision, and single-turn versus iterative reinforcement learning—we systematically explore the fundamental components of this unified approach.
    \item Using this unified framework, we investigate the mechanisms underlying effective reinforcement learning and outline several promising research avenues for further advancing reinforcement learning methodologies for large language models (LLMs).
\end{itemize}

\subsection{Summary of Evaluations and Metrics}

\begin{itemize}[topsep=0pt]
    \item \textbf{English and Chinese Mathematical Reasoning}: We undertake detailed evaluations of our models on both English and Chinese mathematical benchmarks, spanning a range of problems from elementary to collegiate levels. The English benchmarks assessed include GSM8K \citep{gsm8k}, MATH \citep{MATH}, SAT \citep{llemma}, OCW Courses \citep{minerva}, and MMLU-STEM \citep{mmlu}. For Chinese, the selected benchmarks are MGSM-zh \citep{mgsm}, CMATH \citep{wei2023cmath}, Gaokao-MathCloze \citep{agieval}, and Gaokao-MathQA \citep{agieval}. The evaluation protocol examines both the capability to generate fully self-explanatory textual solutions (without reliance on external tools) and the ability to utilize Python for problem-solving.

    On the English benchmark suites, DeepSeekMath-Base performs comparably to the proprietary Minerva 540B \citep{minerva} and outperforms all open-source base models (including Mistral 7B \citep{mistral} and Llemma-34B \citep{llemma}), regardless of their pre-training on mathematical data, often by substantial margins. Particularly for Chinese tasks, DeepSeekMath-Base exhibits marked superiority, which may be attributed to our inclusive pre-training strategy that incorporates high-quality mathematical data in multiple languages, instead of restricting to English-only sources as done in earlier studies \citep{minerva,llemma}. Following instruction-based tuning and reinforcement learning, DeepSeekMath-Instruct and DeepSeekMath-RL achieve state-of-the-art results—most notably, surpassing 50\% accuracy on the highly challenging, competition-level MATH dataset for the first time among open-source models.
\end{itemize}

\item \textbf{Formal Mathematics}: We assess DeepSeekMath-Base's capabilities on the informal-to-formal theorem proving benchmark introduced by \citep{dsp_proof}, utilizing the miniF2F dataset \citep{minif2f} with Isabelle \citep{isabelle} as the proof assistant. DeepSeekMath-Base attains strong results in few-shot autoformalization scenarios.
\item \textbf{Natural Language Understanding, Reasoning, and Code}: To provide a holistic evaluation of the model’s abilities in general comprehension, reasoning, and programming, we apply DeepSeekMath-Base to several benchmarks: Massive Multitask Language Understanding (MMLU) \citep{mmlu}, featuring 57 varied multiple-choice domains; BIG-Bench Hard (BBH) \citep{bbh}, which involves 23 tasks that emphasize multi-step reasoning; along with HumanEval \citep{codex} and MBPP \citep{mbpp}, which are widely recognized coding model benchmarks. Our findings indicate that mathematical pre-training enhances both language comprehension and reasoning outcomes.

\section{Math Pre-Training}

\subsection{Data Collection and Decontamination} This section describes the methodology for curating the DeepSeekMath Corpus from Common Crawl data. Illustrated in Figure \ref{fig:data_collect}, we present an iterative procedure for large-scale acquisition of mathematical web documents. Beginning with a seed dataset—comprised of a small but high-quality sample of mathematical material—the process generalizes to other specialized domains like computer code.

Our initial seed corpus is OpenWebMath \citep{openwebmath}, a rigorously vetted selection of mathematical content from the web. Leveraging this, we train a fastText classifier \citep{joulin2016fasttext} to retrieve further math-relevant web content. For training, we draw 500,000 entries from OpenWebMath to serve as positive samples, paired with an equal number of negatives randomly taken from Common Crawl. The open-source implementation\footnote{\url{https://fasttext.cc}} is employed, with hyperparameters set to a 256-dimensional vector space, a learning rate of 0.1, maximum n-gram length of 3, minimum word occurrence of 3, and 3 epochs. Prior to classification, we shrink the Common Crawl dataset via both exact and approximate URL deduplication, yielding a set of 40B HTML pages. The trained fastText model is then used to extract pages resembling the seed set. Subsequently, candidate documents are ranked according to their fastText scores, and only the highest scoring are retained to exclude low-quality math material. The impact of data scale is evaluated by pre-training on sets consisting of the top 40B, 80B, 120B, and 160B tokens. For the initial round, we select the best 40B-token subset.

Following the initial phase of data acquisition, a significant number of mathematical web pages were still missing from our dataset. This is largely due to the limited diversity present in the set of positive examples used for training the fastText model. To address this, we expanded our seed corpus by sourcing additional mathematical web domains, with the aim of enhancing fastText's performance. The process began by partitioning the entirety of Common Crawl into distinct domains, defined as collections of web pages sharing an identical base URL. For each domain, we computed the fraction of pages acquired during the first data collection round. Domains where more than 10\% of pages had been gathered were designated as math-centric (e.g., \url{mathoverflow.net}).

We then undertook manual annotation of URLs pointing to mathematical content within these domains (such as \url{mathoverflow.net/questions}). The remaining, previously uncollected pages associated with these URLs were integrated into the seed set. This strategy enabled us to amass a broader set of positive examples and subsequently retrain the fastText model to more effectively recover mathematical data in further rounds. By the conclusion of four collection iterations, the dataset comprised 35.5M math-related web pages, amounting to a total of 120B tokens. We observed that, in the fourth round, approximately 98\% of the data had already been harvested in the prior iteration, prompting us to halt additional data collection.

In line with efforts to prevent contamination of benchmarks, we applied the filtering procedure described in \cite{deepseek-coder}, systematically excluding web pages containing either questions or answers from established English mathematical benchmarks (e.g., GSM8K \citep{gsm8k} and MATH \citep{MATH}) and Chinese benchmarks (such as CMATH \citep{wei2023cmath} and AGIEval \citep{agieval}). Specifically, we filtered out any text passage with an exact 10-gram match to any substring from the benchmark datasets. For benchmark items under 10 grams but containing at least 3 grams, the filter relied on exact substring matching to ensure removal of potentially contaminated content from the training set.

\subsection{Validating the Quality of the DeepSeekMath~Corpus}

To assess the relative quality of the DeepSeekMath~Corpus, we carried out pre-training experiments, benchmarking it against several recently published math-oriented corpora:

\begin{itemize}[topsep=0pt]
    \item \textbf{MathPile} \citep{mathpile}: This resource amalgamates 8.9B tokens from varied sources, including textbooks, Wikipedia, ProofWiki, CommonCrawl, StackExchange, and arXiv, with more than 85\% originating from arXiv;
    \item \textbf{OpenWebMath} \citep{openwebmath}: A 13.6B-token dataset consisting of CommonCrawl entries screened for mathematical relevance;
    \item \textbf{Proof-Pile-2} \citep{llemma}: A collection merging OpenWebMath, AlgebraicStack (10.3B tokens of math code), and arXiv publications (28.0B tokens). For our experiments on Proof-Pile-2, we adhere to the 2:4:1 ratio of arXiv:Web:Code established in \cite{llemma}.
\end{itemize}

\subsubsection{Training Setting}
\label{sec:quality-policy}
Mathematical training is carried out on a general-purpose language model consisting of 1.3B parameters, architecturally aligned with the DeepSeek LLM family \citep{deepseek-llm}, and referred to as DeepSeek-LLM 1.3B in this context. Each mathematical dataset is used to independently train a model for 150B tokens. The entire set of experiments utilizes the resource-efficient and lightweight HAI-LLM \citep{haillm} training environment. Adhering to the methodologies described for DeepSeek LLMs, we employ the AdamW optimizer \citep{adamW}, setting $\beta_1=0.9$, $\beta_2=0.95$, and $\mathrm{weight\_decay}=0.1$. The learning rate schedule adopts a multi-step strategy, with an initial warmup phase lasting 2,000 steps, reaching its maximum value, then reducing to 31.6\% at the 80\% mark of training completion, and declining further to 10.0\% of the peak by 90\% of total training. The maximal learning rate is set at 5.3e-4. Training is performed with a batch size of 4M tokens, utilizing a 4K token context window.

\subsubsection{Evaluation Results}

\textbf{DeepSeekMath~Corpus exhibits exceptional quality, encompasses diverse languages in mathematical material, and significantly surpasses others in scale.}

\begin{itemize}[topsep=0pt]
    \item \textbf{High-quality}:
    Downstream model effectiveness is measured on eight mathematical evaluation sets via few-shot chain-of-thought prompting \cite{cot}. According to Table \ref{tab:corpora_comparison}, models trained on DeepSeekMath~Corpus distinctly outperform counterparts. Furthermore, Figure \ref{fig:corpora_comparisons} demonstrates that DeepSeekMath-trained models outperform Proof-Pile-2 at 50B tokens (representing a full epoch over Proof-Pile-2), reinforcing the superior average quality of DeepSeekMath~Corpus. 
    \item \textbf{Multilingual}:
    Data within DeepSeekMath~Corpus is distributed across several languages, with English and Chinese forming the largest segments. Table \ref{tab:corpora_comparison} indicates that training with DeepSeekMath~Corpus significantly improves mathematical reasoning capabilities in both English and Chinese. By comparison, extant mathematical datasets—mainly English-focused—either show negligible gains or can even impair performance on Chinese mathematical reasoning tasks.
    \item \textbf{Large-scale}:
    The scale of DeepSeekMath~Corpus far exceeds that of existing mathematical datasets. As visualized in Figure \ref{fig:corpora_comparisons}, DeepSeek-LLM 1.3B models trained with DeepSeekMath~Corpus display accelerated learning progress and achieve persistent improvements. Meanwhile, baseline corpora, which are notably smaller and utilized for multiple epochs, lead to rapid saturation in model performance.
\end{itemize}

\subsection{Training and Evaluating DeepSeekMath-Base 7B}

This section presents DeepSeekMath-Base 7B, a foundational model designed for robust mathematical reasoning. The model builds upon DeepSeek-Coder-Base-v1.5 7B \citep{deepseek-coder} and is trained on 500 billion tokens. The data distribution used during training consists of 56\% DeepSeekMath Corpus, 4\% AlgebraicStack, 10\% arXiv, 20\% code from Github, and 10\% natural language samples from Common Crawl in both English and Chinese. Training primarily adheres to the configuration outlined in Section \ref{sec:quality-policy}, with two modifications: the learning rate peaks at 4.2e-4 and the batch size is set to 10 million tokens.

A thorough evaluation of DeepSeekMath-Base 7B is conducted, targeting its proficiency in independently generating comprehensive mathematical solutions, utilizing tools to resolve mathematical questions, and performing formal theorem verification. Additionally, we extend the evaluation beyond mathematical tasks to encompass the model’s general capabilities in language understanding, reasoning, and programming.

\paragraph{Mathematical Problem Solving with Step-by-Step Reasoning}

DeepSeekMath-Base is systematically assessed on its ability to solve mathematical questions through few-shot chain-of-thought prompting \citep{cot} over eight benchmark datasets in both English and Chinese. These benchmarks include tasks requiring quantitative reasoning (such as GSM8K \citep{gsm8k}, MATH \citep{MATH}, and CMATH \citep{wei2023cmath}) as well as multiple-choice assessments (including MMLU-STEM \citep{mmlu} and Gaokao-MathQA \citep{agieval}), thereby spanning a wide array of mathematical disciplines from elementary through collegiate levels.

As indicated in Table \ref{tab:base_math_cot}, DeepSeekMath-Base 7B consistently achieves the highest scores across all eight benchmarks compared to open-source base models, which include prominent models such as Mistral 7B \citep{mistral} and the recently introduced Llemma 34B \citep{llemma}, the latter being specifically trained on the Proof-Pile-2 corpus \citep{llemma}. Of particular note is its performance on the competition-caliber MATH dataset, where DeepSeekMath-Base exceeds the leading open-source base models by more than 10 percentage points, and even surpasses Minerva 540B \citep{minerva}—a proprietary model of substantially larger scale (77 times the parameters), built upon PaLM \citep{palm} and subject to further fine-tuning on mathematical material.

\paragraph{Mathematical Problem Solving with Tool Use} 
For program-assisted mathematical reasoning, we evaluate using GSM8K and MATH, employing few-shot program-of-thought prompting \citep{pot,pal}. Models are prompted to construct a Python script to address each problem, with access to computational libraries such as \textit{math} and \textit{sympy} to facilitate sophisticated operations. The output generated from running these scripts is taken as the model’s final answer. Table \ref{tab:base_math_pot_proof} illustrates that DeepSeekMath-Base 7B establishes a new state-of-the-art among open-source base models, outperforming Llemma 34B in this setting.

\paragraph{Formal Mathematics}

Automated formal proof generation plays an essential role in guaranteeing the correctness and robustness of mathematical proofs, a domain garnering increasing interest. We assess DeepSeekMath-Base 7B on informal-to-formal proof tasks as described in \citep{dsp_proof}: given an informal problem statement, its formal equivalent, and an informal proof, the model is asked to produce a formal proof. Evaluation is conducted using miniF2F \citep{minif2f}, a benchmark targeting formalization of Olympiad-level mathematical problems. For each instance, we prompt the model to generate a formal proof in Isabelle using a few-shot approach. Adopting the methodology in \cite{dsp_proof}, we have models produce proof sketches, then utilize the Sledgehammer automated prover \citep{sledgehammer} to complete any unfinished reasoning. According to Table \ref{tab:base_math_pot_proof}, DeepSeekMath-Base 7B achieves notably strong results in automatic proof formalization.

\paragraph{Natural Language Understanding, Reasoning, and Code}

To measure model capacity for language comprehension, we test on MMLU \citep{mmlu}; for reasoning ability, we use BBH \citep{bbh}; and coding skill is evaluated via HumanEval \citep{codex} and MBPP \citep{mbpp}. As reported in Table \ref{tab:understanding_reasoning_code}, DeepSeekMath-Base 7B shows substantial gains on MMLU and BBH in comparison with its predecessor, DeepSeek-Coder-Base-v1.5 \citep{deepseek-coder}, underscoring the advantageous effects of mathematical training on broader language understanding and logical reasoning. The inclusion of code-related tokens in its continued training allows DeepSeekMath-Base 7B to maintain comparable performance to DeepSeek-Coder-Base-v1.5 on both coding benchmarks. Overall, DeepSeekMath-Base 7B markedly surpasses the generalist Mistral 7B \citep{mistral} on all three evaluated reasoning and coding tasks.

\section{Supervised Fine-Tuning}

\subsection{SFT Data Curation}

We compile an instruction-tuning dataset for mathematics that encompasses both English and Chinese tasks, covering a broad spectrum of mathematical domains and a range of difficulties. Each problem is accompanied by an annotated solution in chain-of-thought (CoT) \citep{cot}, program-of-thought (PoT) \citep{pot,pal}, or tool-integrated reasoning formats \citep{tora}. The finalized training set comprises 776K examples.

\begin{itemize}[topsep=0pt]
    \item \textbf{English mathematical datasets}: GSM8K and MATH are enriched with tool-integrated solution annotations. We also include a selected portion of MathInstruct \citep{MathInstruct}, and examples from the Lila-OOD training corpus \citep{lila} where responses are generated using CoT or PoT approaches. This collection spans diverse mathematical subfields, such as algebra, probability, number theory, calculus, and geometry.
    \item \textbf{Chinese mathematical datasets}: Our dataset incorporates K-12 math problems written in Chinese, distributed across 76 distinct topics (e.g., linear equations). Solutions are provided using both CoT and tool-integrated methodologies.
\end{itemize}

\subsection{Training and Evaluating DeepSeekMath-Instruct 7B}

Here we describe the training process for DeepSeekMath-Instruct 7B, which is fine-tuned for mathematical tasks based on DeepSeekMath-Base. During preparation, training samples are concatenated in a random order up to a 4K token context window. The model is optimized over 500 steps with a batch size set to 256, utilizing a constant learning rate of 5e-5.

To assess the mathematical capability of the models, both with and without tool assistance, we conduct evaluations on four quantitative reasoning benchmarks, spanning English and Chinese. Our model's results are compared to state-of-the-art contemporaneous systems, including:

\begin{itemize}[topsep=0pt]
    \item \textbf{Closed-source models}: This set features (1) the GPT lineup, particularly GPT-4 \citep{gpt4} and GPT-4 Code Interpreter~\footnote{\url{https://openai.com/blog/chatgpt-plugins\#\#code-interpreter}}, (2) Gemini Ultra and Pro \citep{gemini}, (3) Inflection-2 \citep{inflection-2}, (4) Grok-1~\footnote{\url{https://x.ai/model-card}}, as well as leading releases from Chinese developers, such as (5) Baichuan-3~\footnote{\url{https://www.baichuan-ai.com}}, and (6) the current GLM-4~\footnote{\url{https://open.bigmodel.cn/dev/api\#glm-4}} from the GLM series \citep{glm}.
\end{itemize}

These models are designed for general applications, the majority of which have undergone alignment processes. 
    \item \textbf{Open-source models} considered in our study span: general-purpose models such as (1) DeepSeek-LLM-Chat 67B \citep{deepseek-llm}, (2) Qwen 72B \citep{qwen}, (3) SeaLLM-v2 7B \citep{seallm}, and (4) ChatGLM3 6B \citep{chatglm3}, as well as models with mathematical performance enhancements, including (5) InternLM2-Math 20B~\footnote{\url{https://github.com/InternLM/InternLM-Math}}, which extends InternLM2 through additional math-specific training and subsequent instruction tuning; (6) Math-Shepherd-Mistral 7B, which utilizes PPO training \citep{schulman2017proximal} applied to Mistral 7B \citep{mistral} with a process-supervised reward model; (7) the WizardMath suite \citep{wizardmath}, which augments the mathematical reasoning capabilities of Mistral 7B and Llama-2 70B \citep{llama2} using evolve-instruct—a technique leveraging AI-generated instructions for instruction tuning—and PPO, with training data predominantly from GSM8K and MATH; (8) MetaMath 70B \citep{metamath}, which is a fine-tuned version of Llama-2 70B using an expanded GSM8K and MATH dataset; (9) ToRA 34B \cite{tora}, which adapts CodeLlama 34B via fine-tuning for enhanced mathematical reasoning with integrated tool usage; and (10) MAmmoTH 70B \citep{MathInstruct}, an instruction-tuned Llama-2 70B model trained on MathInstruct.
\end{itemize}

Referencing Table \ref{tab:sft_rl_math}, within the evaluation protocol prohibiting tool assistance, DeepSeekMath-Instruct 7B attains superior performance in stepwise mathematical reasoning. Of particular note, on the MATH benchmark—geared toward competition-level problems—our model outperforms all other open-source models as well as most proprietary solutions (such as Inflection-2 and Gemini Pro) by a margin of at least 9\% in absolute terms. This holds true even in comparison to models with much larger parameter counts (e.g., Qwen 72B) or those optimized using math-oriented reinforcement learning approaches (e.g., WizardMath-v1.1 7B). DeepSeekMath-Instruct performs comparably to leading Chinese proprietary systems, such as GLM-4 and Baichuan-3, on MATH, although it still falls short of GPT-4 and Gemini Ultra.

When shifting to an evaluation regime that permits the combination of natural language reasoning and programmatic tool use, DeepSeekMath-Instruct 7B reaches almost 60\% accuracy on MATH, overtaking all prior open-source models. For other evaluation sets, our model demonstrates performance on par with DeepSeek-LLM-Chat 67B, the previously leading model despite its tenfold larger scale.

\section{Reinforcement Learning}

\subsection{Group Relative Policy Optimization}
Reinforcement learning (RL) has established its efficacy in further refining the mathematical reasoning skills of LLMs subsequent to the Supervised Fine-Tuning (SFT) phase \citep{wang2023math, wizardmath}. In this section, we introduce our proposed reinforcement learning method, Group Relative Policy Optimization (GRPO), which is both efficient and effective.

\subsubsection{From PPO to GRPO}
Proximal Policy Optimization (PPO) \citep{schulman2017proximal} represents a popular actor-critic RL method widely applied during the RL-based fine-tuning of LLMs \citep{ouyang2022training}. PPO updates model parameters by maximizing the following clipped surrogate objective:

\begin{equation}
\footnotesize
    \mathcal{J}_{PPO}(\theta) = \mathbb{E}{[q \sim P(Q), o \sim \pi_{\theta_{old}}(O|q)]} \frac{1}{|o|} \sum_{t=1}^{|o|} \min \left[ \frac{\pi_\theta(o_{t} | q, o_{<t})}{\pi_{\theta_{old}}(o_{t} | q, o_{<t})} A_{t}, \text{clip} \left( \frac{\pi_\theta(o_{t} | q, o_{<t

In this context, $\pi_{\theta}$ and $\pi_{\theta_{old}}$ represent the current and previous policy models, respectively, while $q$ and $o$ denote questions and outputs, which are drawn from the question dataset and the old policy $\pi_{\theta_{old}}$. The parameter $\varepsilon$ acts as a clipping hyper-parameter utilized in PPO to promote training stability. The advantage estimate, $A_t$, is derived via Generalized Advantage Estimation (GAE) \citep{gae}, leveraging the rewards $\{r_{\ge t}\}$ along with an associated learned value function $V_{\psi}$. Consequently, PPO mandates concurrent training of a value function alongside the policy model, and to curb excessive reward model optimization, it is standard to include a per-token KL divergence penalty relative to a reference model in the token-wise reward signal, as described in \citep{ouyang2022training}:

\begin{equation}
    r_{t} = r_\varphi(q, o_{\le t}) - \beta \log\frac{\pi_{\theta}(o_{t}|q, o_{<t})}{\pi_{ref}(o_{t}|q, o_{<t})},
\label{eq:PPO-reward}
\end{equation}

Here, $r_\varphi$ denotes the reward model, $\pi_{ref}$ refers to the reference model—commonly initialized from the original SFT model—and $\beta$ serves as the scaling factor for the KL penalty.

Given that the value function in PPO is typically implemented as a model roughly equal in scale to the policy model, this incurs significant overhead in both computation and memory. Furthermore, the value function is used in RL as a baseline to reduce variance during advantage calculation. In the large language model (LLM) setting, however, the reward model often provides feedback only on the last token, posing challenges for training a value function that is precise across all tokens. To remedy this issue, as depicted in Figure \ref{fig:grpo}, we introduce Group Relative Policy Optimization (GRPO), which eliminates the necessity for value function approximation required in PPO. Instead, GRPO uses the mean reward of several sampled responses, generated for the same question, as a baseline. More concretely, for each question $q$, GRPO samples a set $\{o_1, o_2, \cdots, o_G\}$ of outputs from the old policy $\pi_{\theta_{old}}$, and updates the policy model to maximize:

\begin{equation}
\footnotesize
\begin{split}
    \mathcal{J}_{GRPO}(\theta) &= \mathbb{E}{[q \sim P(Q), \{o_i\}_{i=1}^G \sim \pi_{\theta_{old}}(O|q)]}  \\
    & \frac{1}{G}\sum_{i=1}^G\frac{1}{|o_i|} \sum_{t=1}^{|o_i|} \left\{ \min \left[ \frac{\pi_\theta(o_{i,t} | q, o_{i,<t})}{\pi_{\theta_{old}}(o_{i,t} | q, o_{i,<t})} \hat{A}_{i,t}, \text{clip} \left( \frac{\pi_\theta(o_{i,t} | q, o_{i,<t})}{\pi_{\theta_{old}}(o_{i,t} | q, o_{i,<t})}, 1 - \varepsilon, 1 + \varepsilon \right)  \hat{A}_{i,t} \right] - \beta \mathbb{D}_{KL}\left[\pi_{\theta} || \pi_{ref}\right]\right\} ,
\end{split}
\label{eq:GRPO-obj}
\end{equation}

In this expression, $\varepsilon$ and $\beta$ remain as hyper-parameters, and $\hat{A}_{i,t}$ is the advantage computed using only the relative rewards within each output group, a method explained in detail in subsequent subsections. The group-relative computation of advantages in GRPO naturally fits with reward models, since these models are often developed from comparison datasets based on the same input question. Also of note, in contrast to applying the KL penalty inside the reward as in standard PPO, GRPO incorporates the KL divergence term directly into the loss, thus avoiding added complexity in the calculation of $\hat{A}_{i,t}$. Additionally, unlike the KL penalty formulation used in (\ref{eq:PPO-reward}), we estimate the KL divergence using an unbiased estimator as in \citep{kl_approx}:

\begin{equation

\begin{algorithm}[t]
  \small
  \caption{Iterative Group Relative Policy Optimization}
  \textbf{Input} Initial policy model $\pi_{\theta_{\text{init}}}$; reward models $r_\varphi$; task prompts $\mathcal{D}$; 
  hyperparameters $\varepsilon$, $\beta$, $\mu$
  \begin{algorithmic}[1]
    \State Initialize policy model $\pi_\theta \leftarrow \pi_{\theta_{\text{init}}}$
    \For{iteration = 1, \dots, I}
       \State Set reference model $\pi_{ref} \leftarrow \pi_{\theta}$
      \For{step = 1, \dots, M}
      \State Randomly select a mini-batch $\mathcal{D}_b$ from $\mathcal{D}$
      \State Assign current policy to old policy: $\pi_{\theta_{old}} \leftarrow \pi_{\theta}$ 
      \State For every query $q \in \mathcal{D}_b$, generate $G$ outputs $\{o_i\}_{i=1}^G \sim \pi_{\theta_{old}} (\cdot \mid q) $
      \State For each output $o_i$, obtain reward values $\{r_i\}_{i=1}^{G}$ using $r_{\varphi}$
      \State Compute the group-relative advantage $\hat{A}_{i,t}$ at token $t$ for each $o_i$
      \For{GRPO iteration = 1, \dots, $\mu$}
        \State Update $\pi_{\theta}$ by optimizing the GRPO objective (Equation \ref{eq:GC-GRPO})
      \EndFor
    \EndFor \State Improve $r_\varphi$ incrementally using a replay strategy for continual training.
    \EndFor 
  \end{algorithmic}
  \textbf{Output} $\pi_\theta$
  \label{alg:iter-grpo}
\end{algorithm}

\subsubsection{Outcome Supervision RL with GRPO} In this setting, for each prompt $q$, a batch of candidate outputs $\{o_1, o_2, \cdots, o_G\}$ is sampled from the previous policy $\pi_{\theta_{old}}$. A reward model subsequently assigns a score to each output, producing a set of $G$ reward values $\mathbf{r} = \{r_1, r_2, \cdots, r_G\}$. These rewards are then standardized: the mean of the group is subtracted and the result is divided by the group’s standard deviation. Under outcome supervision, the resulting normalized reward is used as the advantage for all tokens within $o_i$, i.e., $\hat{A}_{i, t} = \widetilde{r}_i = \frac{r_i- {\rm mean}(\mathbf{r})}{{\rm std}(\mathbf{r})}$, where $\widetilde{r}_i$ denotes the normalized score for $o_i$. The policy is then updated to maximize the objective provided in equation (\ref{eq:GRPO-obj}).

\subsubsection{Process Supervision RL with GRPO} Since outcome-level feedback restricts supervision to only the end of an output, it can be insufficient for guiding learning in tasks that require detailed or stepwise reasoning, such as mathematics. Building on \cite{wang2023math}, process supervision augments the learning signal by supplying a reward at the completion of every reasoning step. More formally, given a query $q$ and $G$ sampled solutions $\{o_1, o_2, \cdots, o_G\}$, a dedicated process reward model evaluates each reasoning step in every output, resulting in the collection of step-wise reward sets: $\mathbf{R} = \{ \{r_1^{index(1)},\cdots,r_1^{index(K_1)}\}, \cdots,  \{r_G^{index(1)},\cdots,r_G^{index(K_G)}\} \}$, with $index(j)$ marking the final token of the $j$-th step, and $K_i$ representing the number of steps in output $i$. These step-level rewards are also standardized as $\widetilde{r}_i^{index(j)} = \frac{r_i^{index(j)} - {\rm mean(\mathbf{R})}}{{\rm std(\mathbf{R})}}$. For each token $t$, the advantage is determined by summing all normalized step rewards for steps ending at or after $t$: $\hat{A}_{i, t} = \sum_{index(j) \ge t

\subsubsection{Iterative RL with GRPO} During the progression of reinforcement learning, the previously trained reward model may no longer provide adequate supervision for the evolving policy model. To address this, we investigate iterative RL with GRPO. As described in Algorithm \ref{alg:iter-grpo}, this approach involves generating fresh datasets for the reward model utilizing outputs sampled from the current policy model. The reward model is incrementally updated via a replay buffer, incorporating 10\% of prior data to preserve historical knowledge. Subsequently, the policy model is assigned as the new reference model and further optimized against the updated reward model in a continual manner.

\subsection{Training and Evaluating DeepSeekMath-RL}

Our reinforcement learning experiments leverage DeepSeekMath-Instruct 7B as the foundation. The RL training data is composed exclusively of chain-of-thought style questions from GSM8K and MATH present within the SFT dataset, yielding approximately 144K instances. We purposefully omit other SFT-derived questions to systematically assess RL's influence on benchmarks that remain unseen during the RL process. The construction of reward model training datasets aligns with the methodology detailed in \citep{wang2023math}. The initial reward model is trained on DeepSeekMath-Base 7B, with a learning rate set to 2e-5. For GRPO, we utilize a policy learning rate of 1e-6 and specify a KL divergence coefficient of 0.04. Each question is associated with $64$ generated responses. The maximum response length allowed is 1024 tokens, and both the batch and training sizes are 1024. The policy model undergoes a single update after each exploration round. Benchmark evaluation of DeepSeekMath-RL 7B follows the same procedure as that of DeepSeekMath-Instruct 7B. In the context of DeepSeekMath-RL 7B, GSM8K and MATH with chain-of-thought reasoning constitute in-domain benchmarks, whereas all other tasks are treated as out-of-domain.

Table \ref{tab:sft_rl_math} presents comparative results for open- and closed-source models employing both chain-of-thought and tool-augmented reasoning strategies across English and Chinese evaluation sets. Our analysis reveals the following: 1) DeepSeekMath-RL 7B achieves accuracy rates of 88.2\% on GSM8K and 51.7\% on MATH when using chain-of-thought reasoning. These results exceed those of all open-source models ranging from 7B to 70B parameters, and also outperform most closed-source alternatives. 2) Importantly, DeepSeekMath-RL 7B is exclusively trained on instruction-tuning data in the chain-of-thought format from GSM8K and MATH, beginning from the DeepSeekMath-Instruct 7B baseline. Despite this limited training corpus, it consistently outperforms DeepSeekMath-Instruct 7B in all measured aspects, underlining the impact of reinforcement learning.

\section{Discussion}
This section outlines the insights gained from our pre-training and reinforcement learning studies.

\subsection{Insights from Pre-Training}

We begin by describing observations from pre-training experiments. Unless indicated otherwise, the methodology mirrors the protocol specified in Section \ref{sec:quality-policy}. For this part, references to the DeepSeekMath Corpus pertain to the 89B-token dataset assembled during the second round of data gathering.

\subsubsection{The Influence of Code Training on Mathematical Reasoning}

While widely assumed but insufficiently substantiated, the belief persists that code training enhances reasoning skills. Our findings provide empirical support for this notion within the context of mathematics: incorporating code training strengthens models' capabilities for mathematical reasoning, regardless of the use of tools.

To examine the effects of code training on mathematical reasoning, we explored two main training paradigms: a two-stage and a single-stage training process:

\textbf{Two-Stage Training}

\begin{itemize}[topsep=0pt]
    \item \textbf{Code Training for 400B Tokens $\rightarrow$ Math Training for 150B Tokens}: DeepSeek-LLM 1.3B is first pretrained on 400B code tokens, followed by continued training on 150B tokens of math data.
    \item \textbf{General Training for 400B Tokens $\rightarrow$ Math Training for 150B Tokens}: For comparison, we replace code tokens with general domain tokens—sampled from DeepSeek-AI’s comprehensive general text dataset—during the initial 400B-token training phase, assessing whether code-specific data provides unique benefits to mathematical reasoning, as opposed to broad-domain text.
\end{itemize}

\textbf{One-Stage Training}

\begin{itemize}[topsep=0pt]
    \item \textbf{Math Training for 150B Tokens}: The model is directly trained for 150B math tokens without pretraining on code or general text.
    \item \textbf{Training on a mixture of 400B Code Tokens and 150B Math Tokens}: Because sequential math training after code pretraining has been observed to diminish coding performance, we explore whether joint training—interleaving code and math tokens in a single stage—can promote both mathematical reasoning and code proficiency while countering catastrophic forgetting.
\end{itemize}

\paragraph{Results}
As illustrated in Table \ref{tab:code-math} and Table \ref{tab:code-math-general-eval}, model performance across several downstream tasks is reported under these varied pretraining regimens.

Introducing code tokens during training leads to measurable improvements in program-assisted mathematical reasoning, in both two-stage and single-stage paradigms. In particular, as indicated by Table \ref{tab:code-math}, pretraining solely on code during the two-stage approach notably strengthens the model's ability to tackle problems in datasets such as GSM8K and MATH, particularly when using Python-based solutions. Subsequent specialized math training further augments these capabilities. For the one-stage setting, combining code and math data within a unified training mix alleviates catastrophic forgetting, as commonly seen after sequential finetuning, and concurrently boosts both coding (see Table \ref{tab:code-math-general-eval}) and program-facilitated mathematical performance (see Table \ref{tab:code-math}).

Beyond program-aided scenarios, code pretraining also facilitates improvements in mathematical reasoning tasks absent tool use. Within two-stage training, commencing with a code pretraining phase confers a moderate advantage, which enhances the impact of the subsequent math-focused training, ultimately achieving top-tier results. Conversely, one-stage blended training with both code and math data is less effective for pure mathematical reasoning (without programming). This may be attributed to the limited representational capacity of DeepSeek-LLM 1.3B, making it difficult for the model to jointly internalize both types of complex data within a single training phase.

\subsubsection{ArXiv Papers Show Limited Effectiveness for Enhancing Mathematical Reasoning}

While arXiv papers are frequently integrated into the mathematical pre-training datasets of language models \citep{minerva,gpt-f,llemma,mathpile}, there remains a lack of comprehensive investigations into their contribution toward mathematical reasoning. Surprisingly, our empirical findings indicate that arXiv papers do not substantially advance mathematical reasoning ability. To assess this, we evaluated models across a spectrum of capacities, notably DeepSeek-LLM 1.3B and DeepSeek-Coder-Base-v1.5 7B \citep{deepseek-coder}, using arXiv-based corpora that underwent distinct preprocessing pipelines:

\begin{itemize}[topsep=0pt]
    \item \textbf{MathPile} \citep{mathpile}: an 8.9B-token dataset assembled using heuristic-based cleaning and filtering; over 85\% of its content derives from scientific arXiv submissions;
    \item \textbf{ArXiv-RedPajama} \citep{redpajama}: the complete collection of arXiv LaTeX source files, curated by excising preambles, comments, macro definitions, and references, encompassing 28.0B tokens.
\end{itemize}

Within our study, DeepSeek-LLM 1.3B was trained on each arXiv corpus for 150B tokens, while DeepSeek-Coder-Base-v1.5 7B was trained for 40B tokens. Across these training regimes, reliance solely on arXiv datasets led to either stagnation or declines in performance over an array of mathematical benchmarks of varying difficulty. These assessments spanned quantitative reasoning tests like GSM8K and MATH (see Table \ref{tab:arxiv-cot}), multiple-choice evaluations such as MMLU-STEM (Table \ref{tab:arxiv-cot}), as well as benchmarks in formal mathematics exemplified by miniF2F (Table \ref{tab:arxiv_atp}).

Nonetheless, there are notable caveats to these findings. Our analysis has yet to address:

\begin{itemize}[topsep=0pt]
    \item The possible influence of arXiv-derived tokens on specialized mathematical tasks beyond the present scope, including tasks like the informalization of theorems, i.e., translating formal mathematical expressions into informal prose;
    \item The effect that arXiv tokens may exert when co-mingled with diverse sources of training data;
    \item The prospect that the advantages of arXiv papers might emerge when leveraging much larger model architectures.
\end{itemize}

Consequently, additional inquiry into these aspects is warranted and remains an avenue for subsequent investigation.

\subsection{Reinforcement Learning: Key Insights}
\subsubsection{Toward a Unified Framework} Here, we present a generalized framework for analyzing a variety of training techniques—including SFT, RFT, DPO, PPO, GRPO—and supplement this with experimental examination of different elements within this framework. In general, the parameter gradient for these training strategies can be encapsulated by the following expression:

\begin{equation}
    \nabla_{\theta}\mathcal{J}_{\textcolor{red}{\mathcal{A}}}(\theta) = \mathbb{E}[\underbrace{(q,o) \sim \textcolor{red}{\mathcal{D}}}_{Data \ Source}]\left( \frac{1}{|o|} \sum_{t=1}^{|o|}  \underbrace{GC_{{\mathcal{A}}}(q, o, t, \textcolor{red}{\pi_{{rf}}})}_{Gradient \ Coefficient}  \nabla_{\theta}\log \pi_{\theta}(o_t | q, o_{<t})\right).
\label{eq:objective}
\end{equation}

This formulation is structured around three central elements: 1) \textit{Data Source $\mathcal{D}$}, specifying the corpus from which samples are drawn for training; 2) \textit{Reward Function $\pi_{{rf}}$}, serving as the generator of feedback signals; and 3) \textit{Algorithm $\mathcal{A}$}, which integrates both the data and reward information to compute the gradient coefficient $GC$ that modulates the intensity of positive or negative updates during learning. We examine several prominent training methodologies through this cohesive lens:

\begin{itemize}
    \item  \textbf{Supervised Fine-tuning (SFT)}: This method adapts a pretrained model by further training it on datasets curated by human annotators for SFT objectives.
    \item \textbf{Rejection Sampling Fine-tuning (RFT)}: RFT involves fine-tuning the SFT model using a subset of outputs generated by the SFT model itself, filtered according to response accuracy to original prompts.
    \item \textbf{Direct Preference Optimization (DPO)}: DPO builds on the SFT model by applying fine-tuning over augmented outputs produced by the SFT model, leveraging pairwise loss functions to optimize preferences.
    \item \textbf{Online Rejection Sampling Fine-tuning (Online RFT)}: In contrast to standard RFT, Online RFT begins with the SFT-initialized model and incrementally fine-tunes it using new outputs sampled from the current, up-to-date policy.
    \item \textbf{PPO/GRPO}: PPO/GRPO procedures also commence from the SFT-pretrained policy, but employ reinforcement based on samples drawn in real time from the evolving model.
\end{itemize}

Table \ref{tab:data_policy} provides a summary of the key elements of the discussed methods. For a comprehensive account of the derivation steps, consult Appendix \ref{app:analysis-rl}.

\paragraph{Observation about Data Source} We classify the sources of training data into two distinct types: online sampling and offline sampling. Online sampling pertains to cases where the training data arises from live exploration using the current policy model during training, whereas offline sampling indicates the use of data generated by the preliminary SFT model. The RFT and DPO methods utilize offline data, in contrast to Online RFT and GRPO, which employ online data sampling.

Figure \ref{fig:rl-analysis} presents our findings that Online RFT outperforms RFT by a large margin across both evaluation benchmarks. In the early training phase, Online RFT exhibits similar performance to RFT, but as training proceeds, Online RFT shows a clear superiority, underscoring the benefits of real-time data acquisition. This aligns with our intuition: initially, both the actor and the SFT model yield similar outputs and the sampled data display limited variance. Over time, however, the actor’s samples diverge more notably, making the advantages of online, up-to-date data sampling increasingly prominent.

\paragraph{Observation about Gradient Coefficient} The algorithm transforms reward signals into gradient coefficients, which in turn drive model parameter updates. In our evaluation, the reward function adopts two forms: `Rule' and `Model.' The `Rule' approach assesses the quality of a response based on answer correctness, while the `Model' variant involves training a reward model—using rule-based judgments as its foundation—to assign scores to responses. Examination of Equations \ref{eq:GC-RFT} and \ref{eq:GC-GRPO} reveals a fundamental distinction: GRPO dynamically adapts its gradient coefficient according to the output of the reward model, enabling selective reinforcement and penalization corresponding to the magnitude of rewards. By contrast, Online RFT lacks this flexibility; it does not penalize wrong answers and treats all correct responses equivalently in terms of reinforcement strength.

As illustrated in Figure \ref{fig:rl-analysis}, GRPO achieves better results than online RFT, underscoring the effectiveness of adjusting positive and negative gradient coefficients. Moreover, GRPO+PS outperforms GRPO+OS, which highlights the advantage of implementing finer, step-sensitive gradient coefficient adjustments. We also investigate iterative RL by conducting two successive iterations in our experiments. According to Figure \ref{fig:iter-rl}, this iterative process substantially boosts performance, with the most significant gains observed after the first iteration.

\subsubsection{Why RL Works?}  
This study applies reinforcement learning to a subset of instruction tuning data and observes marked improvements over the instruction-tuned baseline. To better understand the underlying mechanism of RL's effectiveness, we measure both Pass@K and Maj@K accuracies for the Instruct and RL-optimized models on two separate benchmarks. As depicted in Figure \ref{fig:maj-pass}, RL leads to notable enhancement in Maj@K but does not affect Pass@K. These results suggest that reinforcement learning increases overall robustness in the model’s output distribution, specifically by promoting correct responses among the TopK outputs, rather than directly strengthening underlying abilities. \textbf{Thus, the observed benefit appears to arise from improved selection among candidates, rather than fundamental skill enhancement.} In a similar vein, \citep{wang2023making} identified a \textbf{misalignment problem} in SFT models for reasoning tasks, and demonstrated that preference alignment methods \citep{yuan2023rrhf,song2023preference,wang2023making} can yield better reasoning performance.

\subsubsection{How to Achieve More Effective RL?}
\label{sec:effec_rl}
Our findings confirm that reinforcement learning is highly effective for mathematical reasoning tasks. Furthermore, we propose a unified conceptual framework for understanding various representative training methodologies, positioning them all within the broader spectrum of direct or approximate RL approaches. Equation \ref{eq:objective} captures three essential factors: Data Source, Algorithm, and Reward Function, which we discuss as potential directions for future research.

\paragraph{Data Source} The data source serves as the foundational input for all training paradigms. In the context of reinforcement learning, this refers particularly to unlabeled queries and their generated outputs from the policy model. In this work, we limit ourselves to utilizing questions from instruction tuning and outputs sampled through basic nucleus sampling. We believe this restriction may explain why our RL strategy chiefly improves Maj@K. Future research will extend our RL approach to include out-of-distribution question prompts and incorporate \textbf{advanced sampling (decoding) strategies}, such as tree-search-based methods \citep{yao2023tree}. Additionally, \textbf{efficient inference techniques} \citep{xia-etal-2023-speculative,leviathan2023fast,kwon2023efficient,xia2024unlocking}, which directly influence the exploration effectiveness of policy models, constitute a vital area for future exploration.

\paragraph{Algorithms} Algorithms utilize data along with the reward signal to compute gradient coefficients, which subsequently adjust model parameters. According to Equation \ref{eq:objective}, current approaches generally place full \textbf{TRUST} in the feedback provided by the reward function, either amplifying or suppressing the conditional likelihood of specific tokens. Nevertheless, such trust is problematic as reward signals may not always be dependable, particularly when dealing with intricate tasks. For instance, even the rigorously curated PRM800K datasets \citep{lightman2023let}—annotated by expert annotators—exhibit around a 20\% error rate in their labels\footnote{\url{https://github.com/openai/prm800k/issues/12\#issuecomment-1728491852}}. Therefore, we seek to investigate reinforcement learning algorithms that maintain performance even in the presence of unreliable reward signals. We argue that \textbf{WEAK-TO-STRONG} \citep{burns2023weak} alignment techniques hold the potential to fundamentally transform existing learning methodologies.

\paragraph{Reward Function} The reward function provides the essential feedback signal during training. Within RL frameworks, this function is commonly represented by a neural reward model. We identify three primary avenues for advancing reward model research: 1) \textbf{Improving the generalization capacity of the reward model.} The reward model must extend beyond training distributions to accurately assess novel prompts and sophisticated outputs; failure to do so could mean RL serves only to reinforce current LLM distributions without genuine capability improvement; 2) \textbf{Incorporating uncertainty estimation in reward models.} Modeling this uncertainty may serve as a critical interface between imperfect reward functions and weak-to-strong training procedures; 3) \textbf{Streamlining the construction of precise process reward models} that supply detailed and actionable feedback for tasks involving reasoning \citep{lightman2023let,wang2023math}.

\section{Conclusion, Limitation, and Future Work}

This work introduces DeepSeekMath, an open-source model that surpasses other publicly available systems on the competition-level MATH benchmark and approaches the capabilities observed in proprietary models. DeepSeekMath~begins with the DeepSeek-Coder-v1.5 7B checkpoint and is incrementally trained on a dataset comprising 500B tokens, prominently featuring 120B mathematics-related tokens derived from Common Crawl. Through comprehensive ablation experiments, we identify web-based data sources as rich reservoirs for high-quality mathematical content, while arXiv appears to offer less value than anticipated. Furthermore, we propose Group Relative Policy Optimization (GRPO), a modification of Proximal Policy Optimization (PPO), which demonstrably enhances mathematical reasoning proficiency while requiring reduced memory resources. Experimental outcomes confirm the effectiveness of GRPO, even when DeepSeekMath-Instruct 7B already achieves high benchmark scores. Additionally, we present a cohesive framework for conceptualizing various approaches, and highlight several promising research directions for advancing reinforcement learning efficacy.

Despite the robust quantitative reasoning skills demonstrated by DeepSeekMath~on standard benchmarks, its proficiency in geometry and formal theorem-proving trails behind that of closed-source models. In particular, during preliminary assessments, the model struggled with geometric concepts such as triangles and ellipses, suggesting that pre-training and fine-tuning stages may have suffered from data selection imbalances. Furthermore, given the current parameter count, DeepSeekMath~lags behind GPT-4 in few-shot performance; unlike GPT-4, which benefits from few-shot prompts, DeepSeekMath~does not show significant improvement over zero-shot evaluations. Looking forward, our efforts will focus on refining our data selection strategies to assemble a more robust and diverse pre-training dataset. Moreover, we plan to investigate the possible avenues outlined in Section \ref{sec:effec_rl} for advancing reinforcement learning in large language models.